% !TEX root = tkf.tex
% Theory appendix: structural bias of the BP cumulant under
% column-factorised q. Empirical results live in
% Section~\ref{sec:cumulant-structural-bias-results} (frontmatter-tkf.tex).

\section{Theory: structural bias of the BP cumulant under
column-factorised \texorpdfstring{$q$}{q}}
\label{app:cumulant-structural-bias}

We give a precise derivation of the structural bias displayed
empirically in
Section~\ref{sec:cumulant-structural-bias-results}.  The bias is
not a finite-data noise artifact — it is a property of the
column-factorised variational family used by SVI-VarAnc, and
persists at any $q$ whose per-column marginals match the truth.

\paragraph{Setup.}
On a single branch $e$ of length $t_e$, the per-branch WFST is a
chain HMM over the $5$ states $\{S, M, I, D, E\}$ with chain
transition matrix $\mathbf{T}_e[s, s']$ given by the standard
TKF91 / TKF92 closed forms (\texttt{tkf92\_wfst\_T}, eq.\
\ref{eq:tkf92-wfst}).  The branch's contribution to the data
log-likelihood is $\sum_{n=1}^{L+1} \log T_e[\zeta_{n-1}, \zeta_n]$
where $\boldsymbol{\zeta} = (\zeta_0, \dots, \zeta_{L+1})$ is the chain
trajectory across the $L+2$ columns (including boundaries).
Conventionally $\zeta_n \in \{S, M, I, D, E, \mathrm{Ig}\}$ where $\mathrm{Ig}$
denotes a non-emitting "insertion-gap" column where the branch is
in NYI/NYI or $D$/$D$ joint state and contributes nothing to the
chain transitions; chain transitions are indexed by consecutive
non-$\mathrm{Ig}$ columns.

The variational $q$ factorises across columns: $q(\boldsymbol{\zeta}) =
\prod_{n=1}^L q_n(\zeta_n)$ with each $q_n$ a tree-structured
distribution over edge presence/absence carrying the leaf clamps
of column $n$.  Define the per-column WFST-state probabilities at
branch $e$:
\begin{equation}
P_n^{(e)}(s) = q_n(\text{branch state } e \text{ is } s),
\quad s \in \{M, I, D\},
\qquad
P_n^{(e)}(\mathrm{Ig}) = q_n(\text{NYI},\text{NYI}) + q_n(D, D).
\end{equation}

\paragraph{The cumulant formula and what it computes exactly.}
Section~\ref{sec:bw-tkf92} introduces the BP cumulant $\mathbf{W}_e
\in \mathbb{R}^{5 \times 5}$ as the sufficient statistic the
M-step ingests:
\begin{equation}
\label{eq:cumulant-W}
W_e[s, s']
= \sum_{1 \leq M < N \leq L+1}
   P_M^{(e)}(s) \cdot P_N^{(e)}(s') \cdot
   \prod_{k = M+1}^{N-1} P_k^{(e)}(\mathrm{Ig}).
\end{equation}
Under the factorised $q$, $W_e[s, s']$ is the \emph{exact}
expected count of WFST chain transitions $s \to s'$ on branch $e$:
\begin{align*}
\expect_q\bigl[\#\{\text{chain transitions } s \to s'\}\bigr]
&= \sum_{M < N} q\bigl(\zeta_M = s, \zeta_N = s',
                        \zeta_{M+1} = \cdots = \zeta_{N-1} = \mathrm{Ig}\bigr) \\
&= \sum_{M < N} P_M(s)\,P_N(s')\,\prod_{k=M+1}^{N-1} P_k(\mathrm{Ig})
\;\;=\;\; W_e[s, s'],
\end{align*}
where the second equality uses the column-factorisation of $q$.
The sum over $(M, N)$ counts each chain transition once because
of the explicit $\mathrm{Ig}$-padding constraint: any single trajectory
realisation $\boldsymbol{\zeta}$ with non-$\mathrm{Ig}$ columns at indices
$n_1 < n_2 < \cdots < n_K$ contributes weight 1 to exactly the
adjacent-pair transitions $(n_0, n_1), (n_1, n_2), \dots,
(n_K, n_{K+1})$ (with $n_0 = 0, n_{K+1} = L+1$); for any other
$(M, N)$ pair, at least one intervening column is non-$\mathrm{Ig}$ and
the indicator product is zero.  This $O(L)$ prefix-sum implementation
agrees with the $O(L^2)$ direct sum to machine precision
($6.7\times 10^{-16}$).

\paragraph{The structural bias.}
The factorisation assumption is \emph{not} satisfied by the truth
posterior $p^*(\boldsymbol{\zeta} \mid \text{data})$, even for TKF91
($\ext = 0$).  Although TKF91's generative model is per-position
i.i.d., the per-branch WFST chain has nonzero transition
probabilities $T_e[D, D]$, $T_e[I, I]$, $T_e[I, D]$ and so on
(closed-form expressions in \eqref{eq:tkf92-wfst} reduce, at
$\ext = 0$, to the standard TKF91 Pair HMM weights).  These
chain memories propagate from data conditioning into a
non-factorised joint posterior:
\begin{equation}
p^*\bigl(\zeta_M = s, \zeta_N = s', \zeta_{M+1..N-1} = \mathrm{Ig}
   \mid \text{data}\bigr)
\;\neq\;
\prod_n p^*\bigl(\zeta_n \mid \text{data}\bigr),
\end{equation}
even though the per-column truth marginal $p^*(\zeta_n \mid
\text{data})$ is exactly recoverable by per-column Felsenstein
BP.  The bias of the cumulant against the truth-posterior chain
expectation is the integrated difference between the factorised-q
product of marginals and the truth-posterior joint:
\begin{equation}
\label{eq:cumulant-bias-correct}
\boxed{
W_e[s, s'] - \expect_{p^*}\!\bigl[\#\{s \to s'\} \bigm| \text{data}\bigr]
= \sum_{M < N}\!\Bigl[ p^*_M(s)\,p^*_N(s')\!
   \prod_{k=M+1}^{N-1}\! p^*_k(\mathrm{Ig})
\;-\; p^*\bigl(\zeta_M=s, \zeta_N=s', \zeta_{M+1..N-1}=\mathrm{Ig}
   \bigm| \text{data}\bigr) \Bigr].
}
\end{equation}
The bracketed term is the connected $(N-M+1)$-way correlation of
the truth posterior over the columns $\{M, M+1, \dots, N\}$; it
vanishes only when the truth posterior factorises across these
columns.  For TKF91, $T_e[D,D] > 0$ and $T_e[I,I] > 0$ at
$\ext = 0$ alone suffice to make the right-hand side of
\eqref{eq:cumulant-bias-correct} nonzero on the $(s, s') = (I, I)$
and $(D, D)$ entries: an opened gap is positively correlated with
continued gap by the chain transition kernel itself, even before
any explicit fragment-extension parameter is invoked.

\paragraph{Entropy reward exacerbates the bias on rare entries.}
The ELBO objective decomposes as $\mathrm{ELBO} = \expect_q[\log p(z,
\text{data}; \theta)] - \expect_q[\log q]$ with the entropy term
$-\expect_q[\log q]$ rewarding diffuse $q$.  At ELBO stationarity,
$q$'s per-column marginals balance the data-fit term against the
entropy term; for tied or weakly-fitted logits, the resulting $q$
is slightly more diffuse than the truth posterior on individual
rare states (i.e., $P_n^{(e)}(I)$ above the per-column truth
marginal $p_n^*(I)$).  Through the cumulant
\eqref{eq:cumulant-W}, even a small per-column inflation
$P_n(s) = p_n^*(s) + \epsilon$ contributes a leading
linear-in-$\epsilon$ bias on $W[s, s']$ with positive coefficient
$\partial W / \partial P_n(s) = \sum_{N > n} P_N^{(e)}(s')
\prod_{k=n+1}^{N-1} P_k^{(e)}(\mathrm{Ig}) > 0$, summed over $L$ source
columns; the cumulant amplifies per-column diffuseness via its
$O(L^2)$ pair structure.

\paragraph{Why the EM converges to a fixed point.}
The EM map $\theta \mapsto \mathrm{MLE}(\mathbf{W}(\theta))$ is
self-consistent at the biased fixed point $\theta_\infty$ when
the cumulant bias on $\theta$ saturates.  Each entry of the cumulant
is bounded uniformly in $\theta$:
$W_e[s, s'] \leq L + 1$ on a single branch (each $(M, N)$ pair
contributes at most $1$, and the geometric $\mathrm{Ig}$-product is
summable),
giving a $\theta$-independent upper bound on the aggregated
sufficient-statistic count.  The exposure denominator
$T_{\text{eff}} = \sum_e t_e$ in the BDI rate-MLE
(\texttt{tkf92\_vbem.py:223}) is itself $\theta$-independent ---
it depends only on the tree branch lengths, not on the inferred
rates --- and is bounded below by the smallest tree branch.
Composing these, the rate-MLE $\hat\lambda(\theta) \propto
\sum_e W[\cdot, I] / T_{\text{eff}}$ (with the precise BDI
quadratic-MLE form in \texttt{tkf92\_vbem.py:m\_step\_indel\_quadratic})
is a continuous, bounded self-map on a compact rate interval $[0,
\bar\lambda]$ with $\bar\lambda \leq (L + 1) / \min_e t_e$
(continuity follows because $W(\theta)$ depends continuously on
$\theta$ via the per-edge TKF92 transition matrix
\eqref{eq:tkf92-wfst}, and the M-step quadratic root is continuous
in its coefficients).  Brouwer's fixed-point theorem applies:
a fixed point $\theta_\infty$ exists.  Consistent with this
boundedness argument, the EM trajectory of
Section~\ref{sec:cumulant-structural-bias-results} reaches a finite
biased fixed point and does not diverge.

\paragraph{Implications for parameter inference.}
The structural bias is a property of any inference scheme that
uses the cumulant tensor $\mathbf{W}$ as a sufficient statistic
under a column-factorised variational family.  It does not affect:
(i) the ELBO value itself (which is a coherent $\log p(\text{data})$
lower bound regardless of what the cumulant looks like), nor
(ii) ancestral indel-presence reconstruction
(Section~\ref{sec:results-varanc-vs-fitch}: VarAnc and Fitch
parsimony are statistically indistinguishable on the unified
Pfam test splits, $|\Delta\text{F1}| < 0.006$ across short / hard
/ xhard), nor
(iii) downstream MSA-conditioned rate estimation via
exact-EM (Maraschino, Section~\ref{sec:results-pfam-cherry}).
It does affect any rate estimate computed from the variational
cumulant, including the SVI-VarAnc refinement of cherry-trained
parameters (Section~\ref{sec:results-cherry-vs-msa}), where
parameters drift toward $\theta_\infty$ rather than toward
truth.  We report both initial cherry-trained and $\theta_\infty$
estimates without correction; the structural bias is the principal
known limitation of variational training in this paper.
